\chapter{Objectifs d'apprentissage}\label{III-objectifs-dapprentissage}

Cette section s'efforce d'expliciter et de hiérarchiser les acquis
d'apprentissages associés au chapitre. Ces objectifs sont organisés en
paliers~:

($\mathord{\boldsymbol{\circ}}$) Prérequis ($\mathord{\bullet}$)
Fondamental ($\mathord{\bullet}\mathord{\bullet}$) Standard
($\mathord{\bullet}\mathord{\bullet}\mathord{\bullet}$) Avancé
($\mathord{\bullet}\mathord{\bullet}\mathord{\bullet}\mathord{\bullet}$)
Expert

Sauf mention particulière, la connaissance des démonstrations du
document n'est pas exigible\footnote{l'étude des démonstrations du cours
  peut toutefois contribuer à votre apprentissage, au même titre que la
  résolution d'exercices.}

\runinhead{Lois
conditionnelles}\phantomsection\label{III-lois-conditionnelles}

\begin{itemize}
\tightlist
\item
  $\mathord{\bullet}\mathord{\bullet}$ connaître le théorème de Fubini
  conditionnel
\item
  $\mathord{\bullet}\mathord{\bullet}$ connaître les formules de
  balayage conditionnel
\item
  $\mathord{\bullet}\mathord{\bullet}$ savoir appliquer ces résultats
  pour différents types de loi de probabilité (à densité ou non)
\item
  $\mathord{\bullet}$ connaître le critère d'indépendance qui résulte du
  théorème de Fubini conditionnel
\end{itemize}

\runinhead{Espérance
conditionnelle}\phantomsection\label{III-espuxe9rance-conditionnelle}

\begin{itemize}
\tightlist
\item
  $\mathord{\bullet}$ connaître les deux points de la définition de
  l'espérance conditionnelle
\item
  $\mathord{\bullet}$ connaître les deux points de la définition de
  l'espérance conditionnelle d'une fonction de variables aléatoires
\item
  $\mathord{\bullet}$ connaître et savoir utiliser la formule de
  l'espérance totale
\item
  $\mathord{\bullet}\mathord{\bullet}$ savoir calculer la densité
  conditionnelle d'un certain nombre de composantes sachant les autres
  dans un vecteur gaussien à densité
\end{itemize}

\runinhead{Cas
\(\mathcal{L}^2\)}\phantomsection\label{III-cas-mathcall2}

\begin{itemize}
\tightlist
\item
  $\mathord{\bullet}\mathord{\bullet}$ savoir que la régression linéaire
  est la meilleure approximation linéaire (au sens des moindres carrés)
  d'une variable aléatoire par une autre
\item
  $\mathord{\bullet}\mathord{\bullet}\mathord{\bullet}$ savoir retrouver
  ce résultat
\item
  $\mathord{\bullet}\mathord{\bullet}\mathord{\bullet}$ connaître
  l'interprétation géométrique de l'espérance conditionnelle dans le cas
  \(\mathcal{L}^2\)
\item
  $\mathord{\bullet}$ connaître et savoir utiliser la formule de la
  variance totale
\end{itemize}

\chapter{Lois conditionnelles}\label{III-lois-conditionnelles-1}

\section{Introduction}\label{III-introduction}

On s'est consacré jusqu'à présent à l'étude de variables aléatoires
indépendantes. En pratique cependant, on rencontre souvent des variables
dépendant les unes des autres. Dans le cas de la météo, les variables
température, vitesse du vent et pression en fournissent un exemple. Dans
les approches bayésiennes, on résume l'information disponible sur l'état
du système étudié par la \textbf{loi a priori} et on met à jour notre
connaissance du système en incorporant de l'information supplémentaire
(par exemple des observations). On cherche alors à caractériser la
\textbf{loi a posteriori} de l'état du système, qui est la loi de l'état
sachant l'information supplémentaire. On va ainsi s'attacher dans ce
chapitre à décrire les \textbf{lois conditionnelles} qui vont permettre
de résumer l'information apportée par une variable (ou un vecteur) sur
une autre et s'intéresser en particulier à l'\textbf{espérance
conditionnelle} qui indiquera le comportement moyen d'une variable
conditionnellement à une autre. Ce dernier cas pose le cadre
probabiliste d'un des problèmes fondamentaux en apprentissage
statistique~: l'apprentissage supervisé, où on dispose d'un ensemble de
réalisations d'une variable dont on cherche à prédire le comportement à
partir d'un ensemble de variables dites explicatives (ou prédicteurs).

\section{Lois conditionnelles dans un
couple}\label{III-lois-conditionnelles-dans-un-couple}

Soient deux variables aléatoire \(X\) et \(Y\) définies sur le même
espace probabilisé \((\Omega, \mathcal{A}, \mathbb{P})\). Dans le cas où
\(X\) et \(Y\) sont indépendantes, on a vu que pour tous boréliens
\(B_1\) et \(B_2\) de \(\mathcal{B}(\mathbb{R})\), on a
\[\mathbb{P}(X\in B_1, Y\in B_2)= \mathbb{P}(X\in B_1)\mathbb{P}(Y\in B_2) = \mathbb{P}_X(B_1)\mathbb{P}_Y(B_2) = \int_{B_1}\mathbb{P}_Y(B_2)d\mathbb{P}_X(x),\]
où on a utilisé le
\href{Probabilité\%20II.pdf\%20\#fubiniproba}{théorème de Fubini}.

Du fait de l'indépendance, on a aussi
\(\mathbb{P}_Y(B_2) = \mathbb{P}(Y\in B_2) = \mathbb{P}(Y \in B_2 | X \in B_1) = \mathbb{P}_Y(B_2|X \in B_1)\)
ce qui exprime que pour tout borélien \(B_1\), la loi conditionnelle de
\(Y\) sachant \(X\in B_1\) est identique à la loi de \(Y\).

Lorsque \(X\) et \(Y\) en sont pas indépendantes, on va chercher à
établir une égalité de la forme
\[\mathbb{P}(X\in B_1, Y\in B_2) = \mathbb{P}_X(B_1)\mathbb{P}_Y(B_2 |X\in B_1) = \int_{B_1}\mathbb{P}_{Y|X=x}(B_2)d\mathbb{P}_X(x)\]
et s'intéresser à caractériser la \emph{loi conditionnelle de \(Y\)
sachant \(X=x\)}, que l'on notera donc \(\mathbb{P}_{Y|X=x}\).

De même, pour toute application \(g : \mathbb{R}^2 \to \mathbb{R}\)
mesurable telle que \(g(X,Y)\) admette une espérance (relativement à la
loi du couple \(\mathbb{P}_{X,Y}\)), on voudrait écrire~:
\[\mathbb{E}(g(X,Y)) = \int_{\mathbb{R}} \left( \int_{\mathbb{R}} g(x,y) d\mathbb{P}_{Y|X=x}(y)\right) d\mathbb{P}_X(x)\]

Pour bien fixer les idées, on va décrire spécifiquement les cas où \(X\)
est discrète puis où le couple \((X,Y)\) admet une densité avant
d'aborder le cas général.

\section{Cas où \(X\) est
discrète}\label{III-cas-ouxf9-x-est-discruxe8te}

Dans ce paragraphe, on suppose que la variable aléatoire réelle \(X\)
est discrète, c'est-à-dire que l'ensemble
\(X(\Omega) \subset \mathbb{R}\) des valeurs \(x_k\) prises par \(X\)
est au plus dénombrable.

On peut imposer que \(\forall x \in X(\Omega)\) on ait
\(\mathbb{P}(X=x) > 0\), quitte à modifier \(X\) sur un ensemble de
probabilité nulle. On va ainsi pouvoir utiliser la définition de la
probabilité conditionnelle pour des événements de la forme \(\{X =x\}\).
Ceci permet d'écrire pour tous boréliens \(B_1\) et \(B_2\) de
\(\mathbb{R}\)~: \begin{align*}
\mathbb{P}(X \in B_1, Y \in B_2) &= \sum_{x \in X(\Omega)\cap B_1} \mathbb{P}(X=x, Y\in B_2)\\
                         &= \sum_{x \in X(\Omega)\cap B_1} \mathbb{P}(X=x) \mathbb{P}(Y \in B_2 | X=x)\\
                         &= \int_{B_1} \mathbb{P}(Y \in B_2 | X=x) d\mathbb{P}_X(x)
\end{align*} puisque
\(\mathbb{P}_X = \sum_{x \in X(\Omega)} \mathbb{P}(X=x)\delta_x\). On
obtient ainsi l'écriture souhaitée en posant
\[\mathbb{P}_{Y|X=x}(B_2) = \mathbb{P}(Y \in B_2 | X=x),\,\,\,\forall x \in X(\Omega), \forall B_2\in\mathcal{B}(\mathbb{R}).\]

\addcontentsline{toc}{subsection}{Image de la probabilité
conditionnelle}
\begin{remarque}[{Image de la probabilité
conditionnelle}]\phantomsection\label{III-image-de-la-probabilituxe9-conditionnelle}

\(\mathbb{P}_{Y|X=x}\) ainsi définie est simplement la probabilité sur
\((\mathbb{R},\mathcal{B}(\mathbb{R}))\) image par \(Y\) de la
probabilité conditionnelle \(\mathbb{P}(\cdot|X=x)\) définie sur
\((\Omega,\mathcal{A})\), autrement dit, la \textbf{loi de \(Y\)
relative à \(\mathbb{P}(\cdot|X=x)\)} et non à \(\mathbb{P}\).

La formule ci-dessus s'écrit
\(\mathbb{P}_{X,Y}(B_1 \times B_2) = \int_{B_1} \mathbb{P}(Y \in B_2 | X=x) d\mathbb{P}_X(x)\),
où \(\mathbb{P}_{X,Y}\) est la loi du couple. Elle se généralise à tout
borélien \(B\) de \(\mathbb{R}^2\) de la manière suivante~: \[
\begin{aligned}
\mathbb{P}_{X,Y}(B) &= \mathbb{P}((X,Y)\in B) = \sum_{x \in X(\Omega)} \mathbb{P}(X=x, (x,Y) \in B) \\
      &= \sum_{x \in X(\Omega)} \mathbb{P}(X=x) \mathbb{P}((x,Y) \in B | X=x) \\
      &= \sum_{x \in X(\Omega)} \mathbb{P}(X=x) \mathbb{P}_{Y|X=x}(B_x),
\end{aligned}
\] où \(B_x = \{y\in \mathbb{R}, (x,y) \in B\}\). Ainsi, pour tout \(B\)
borélien de \(\mathbb{R}^2\),

\[\mathbb{E}(1_B(X,Y)) = \int_{\mathbb{R}^2} 1_B(x,y)d\mathbb{P}_{X,Y}(x y) = \int_\mathbb{R}\left(\int_\mathbb{R}1_B(x,y) d\mathbb{P}_{Y|X=x}(y)\right)  d\mathbb{P}_X(x)\]

Par linéarité de l'espérance, on peut ainsi exprimer l'espérance d'une
fonction étagée. Pour avoir le résultat pour une fonction mesurable
positive, on exprime celle-ci comme limite simple d'une suite croissante
de fonctions étagées, et on applique le théorème de convergence
monotone. Enfin, on applique cette construction à \(g_+\) et \(g_-\)
pour une fonction \(g\) de signe quelconque
\(\mathbb{P}_{X,Y}\)-intégrable. En d'autres termes, on reprend le
procédé de construction de l'intégrale de Lebesgue. On obtient ainsi la
formule souhaitée~:
\[\mathbb{E}(g(X,Y)) = \int_{\mathbb{R}} \left( \int_{\mathbb{R}} g(x,y) d\mathbb{P}_{Y|X=x}(y)\right) d\mathbb{P}_X(x).\]
\end{remarque}

\addcontentsline{toc}{subsection}{Pour fixer les idées (1)}
\begin{exemple}[{Pour fixer les idées
(1)}]\phantomsection\label{III-ex1}

Soit \(X \geq 0\) une variable aléatoire à valeurs dans \(\mathbb{N}\)
et \(Y\) une variable aléatoire réelle positive telle que la loi du
couple \(\mathbb{P}_{X,Y}\) vérifie pour tout \(n\in\mathbb{N}\) et tout
borélien \(B_2\) de \(\mathbb{R}\)~:
\[\mathbb{P}_{X,Y} (\{n\}\times B_2) = (1-\alpha)\alpha^n \int_{B_2 \cap \mathbb{R}_+^\ast}e^{-t}\frac{t^n}{n!}dt,\,\,\, 0 < \alpha <1\]
\(\mathbb{P}_{X,Y}\) est bien une probabilité sur \(\mathbb{R}^2\)
puisque par convergence monotone~: \begin{align*}
\mathbb{P}_{X,Y} (\mathbb{R}^2) &= \mathbb{P}_{X,Y}(\mathbb{N}\times \mathbb{R}) \\
                &= \sum_{n\in\mathbb{N}} \mathbb{P}_{X,Y} (\{n\}\times \mathbb{R})\\
                &= \sum_{n\in\mathbb{N}} (1-\alpha)\alpha^n \int_{\mathbb{R}_+^\ast}e^{-t}\frac{t^n}{n!}dt \\
                &= (1-\alpha)\int_{\mathbb{R}_+^\ast}e^{-t} \sum_{n\in\mathbb{N}} \frac{(\alpha t)^n}{n!}dt \\
                &= (1-\alpha)\int_{\mathbb{R}_+^\ast}e^{-(1-\alpha)t} dt = 1
\end{align*}\\
où on aura reconnu la loi exponentielle de paramètre \((1-\alpha)\).
\(\forall n \in \mathbb{N}\),
\[ \int_{\mathbb{R}_+^\ast}e^{-t}\frac{t^n}{n!}dt = \int_{\mathbb{R}_+^\ast}e^{-t}\frac{t^{(n-1)}}{(n-1)!}dt = \ldots = \int_{\mathbb{R}_+^\ast}e^{-t}dt =1\]
par intégrations par parties itérées. La loi marginale de \(X\) s'écrit
donc~:
\[\forall n \in \mathbb{N}, \mathbb{P}(X=n) = \mathbb{P}_{X,Y} (\{n\}\times \mathbb{R}_+^\ast) = (1-\alpha)\alpha^n,\]
loi géométrique de paramètre \((1-\alpha)\). On en déduit la loi
conditionnelle de \(Y\) sachant \(X = n\)~:
\[\mathbb{P}_{Y|X=n}(B_2) = \mathbb{P}(Y \in B_2 | X=n) = \frac{\mathbb{P}_{X,Y} (\{n\}\times B_2)}{\mathbb{P}(X=n)} = \int_{B_2 \cap \mathbb{R}_+^\ast} e^{-t}\frac{t^n}{n!}dt\]
et \(\mathbb{P}_{Y|X=n}\) est la donc la loi gamma de paramètre
\((n+1,1)\).
\end{exemple}

\section{Densités
conditionnelles}\label{III-densituxe9s-conditionnelles}

On suppose maintenant que le couple \((X,Y)\) admet une densité
\(f_{X,Y}\) (par rapport à la mesure de Lebesgue). On note
\(f_X(x) = \int_\mathbb{R}f_{X,Y}(x,y)dy\) (respectivement
\(f_Y(y) = \int_\mathbb{R}f_{X,Y}(x,y)dx\)) la loi marginale de \(X\)
(resp. de \(Y\)). On s'intéresse à caractériser la densité de la
variable \(Y\) connaissant la valeur prise par la variable \(X\), c'est
la \emph{densité conditionnelle} de \(Y\) sachant \(\{X = x\}\)~:

\addcontentsline{toc}{subsection}{Densité conditionnelle}
\begin{proposition}[{Densité
conditionnelle}]\phantomsection\label{III-defdenscond}

La formule suivante définit une densité sur \(\mathbb{R}\), pour tout
\(x \in \mathbb{R}\) tel que \(f_X(x) > 0\).
\[ f_{Y|X=x}(y) = \frac{f_{X,Y}(x,y)}{f_X(x)}.\] Cette fonction
s'appelle la \emph{densité conditionnelle de \(Y\) sachant
\(\{X = x\}\)}. La probabilité conditionnelle de \(Y\) sachant
\(\{X = x\}\) s'écrit ainsi \(\mathbb{P}_{Y|X=x} = f_{Y|X=x}\lambda\),
où \(\lambda\) représente la mesure de Lebesgue.
\end{proposition}

\begin{preuve}
La preuve est immédiate puisque \(f_{Y|X=x}\) est une fonction positive
d'intégrale 1.\hfill$\blacksquare$
\end{preuve}

\runinhead{Exercice -- Dans un triangle
(1) (\(\mathord{\bullet}\))}\phantomsection\addcontentsline{toc}{subsubsection}{Dans
un triangle (1)}\label{III-etb1}
Soient \(X\) et \(Y\) de densité jointe
\(f_{X,Y}(x,y)= \frac{1}{x}1_T (x,y)\) où \(T\) est le triangle
\(T = \{0< y< x < 1\}\).

\begin{enumerate}
\def\labelenumi{\arabic{enumi}.}
\tightlist
\item
  Calculer la densité marginale de \(X\)
\item
  Calculer la densité conditionnelle de \(Y\) sachant \(X=x\).
\end{enumerate}

\solref{III-answer-etb1}

L'interprétation de cette définition est la suivante~: la fonction
\(f_{Y|X=x}\) est la densité de la ``loi conditionnelle de \(Y\) sachant
que \(X = x\)''. Bien sûr, nous avons \(\mathbb{P}(X = x) = 0\) puisque
\(X\) admet une densité, donc la phrase ci-dessus n'a pas réellement de
sens, mais elle se justifie heuristiquement ainsi~: \(dx\) et \(dy\)
étant de ``petits'' accroissements des variables \(x\) et \(y\) et
lorsque \(f\) et \(f_X\) sont continues et strictement positives
respectivement en \((x,y)\) et \(x\)~: \begin{align*}
f_X(x) dx & \approx \mathbb{P}(X \in [x, x+dx])\\
f_{X,Y}(x,y) dx dy & \approx \mathbb{P}(X \in [x, x+dx], Y \in [y, y+dy])\\
\end{align*} Par suite \begin{align*}
f_{Y|X=x} (y) dy & \approx \frac{\mathbb{P}(X \in [x, x+dx], Y \in [y, y+dy])}{\mathbb{P}(X \in [x, x+dx])}\\
                 & \approx \mathbb{P}(Y \in [y, y+dy]|X \in [x, x+dx])\\
\end{align*}

On a alors le résultat suivant qui résout le problème posé en
introduction~:

\begin{proposition}\phantomsection\label{III-proposition}
Pour toute fonction \(g : \mathbb{R}^2 \to \mathbb{R}\) telle que
\(g(X,Y)\) admette une espérance, on a~:
\[\mathbb{E}(g(X,Y)) = \int_\mathbb{R}\left( \int_\mathbb{R}g(x,y)f_{Y|X=x}(y) dy \right) f_X(x) dx,\]
dont on déduit, en prenant \(g=1_{B_1 \times B_2}\), que~:
\[\mathbb{P}(X\in B_1, Y\in B_2) = \int_{B_1} \left( \int_{B_2}f_{Y|X=x}(y) dy \right) f_X(x) dx.\]
\end{proposition}

\begin{preuve}
On a \begin{align*}
\mathbb{E}(g(X,Y)) &= \int_{\mathbb{R}^2} g(x,y) f_{X,Y}(x,y) dy dx \\
             &= \int_{\mathbb{R}^2} g(x,y) f_{Y|X=x}(y)f_X(x) dy dx \\
             &= \int_\mathbb{R}\left( \int_\mathbb{R}g(x,y)f_{Y|X=x}(y) dy \right) f_X(x)dx,
\end{align*} les calculs étant licites par application du
\href{Probabilité\%20II.pdf\%20\#fubiniproba}{théorème de Fubini} et du
fait que l'application
\(x \mapsto \int_\mathbb{R}g(x,y)f_{Y|X=x}(y) dy\) est définie pour
\(f_X(x) >0\), soit presque partout relativement à la mesure
\(\mathbb{P}_X = f_X \lambda\).\hfill$\blacksquare$
\end{preuve}

\section{Cas général}\label{III-cas-guxe9nuxe9ral}

On peut établir le résultat suivant, qui complète le
\href{Probabilité\%20II.pdf\%20\#fubiniproba}{théorème de Fubini} et le
résultat d'existence et d'unicité des mesures produits, et que l'on
admettra.

\addcontentsline{toc}{subsection}{Fubini conditionnel}
\begin{theoreme}[{Fubini
conditionnel}]\phantomsection\label{III-fubinicond}

Soit un couple \((X,Y)\) de variables aléatoires réelles de loi jointe
\(\mathbb{P}_{X,Y}\), il existe une famille
\(\left(\mathbb{P}_{Y|X=x}\right)_{x\in\mathbb{R}}\) de probabilités sur
\((\mathbb{R},\mathcal{B}(\mathbb{R}))\), unique à une égalité
\(\mathbb{P}_X\)-presque partout près\footnote{c'est-à-dire qu'on peut
  définir ces probabilités de la manière qu'on souhaite pour les
  boréliens \(B\) tels que \(\mathbb{P}_X(B)=0\).}, qui vérifie pour
tous \(B_1, B_2\) boréliens de \(\mathbb{R}\)~:
\[ \mathbb{P}_{X,Y}(B_1 \times B_2) = \int_{B_1} \left( \int_{B_2} d\mathbb{P}_{Y|X=x}(y) \right) d\mathbb{P}_X(x).\]
Ces probabilités sont appelées \emph{lois conditionnelles} de \(Y\)
sachant \(X =x\). On a de plus pour toute application
\(g : \mathbb{R}^2 \to \mathbb{R}\) telle que \(g(X,Y)\) admette une
espérance~:
\[\mathbb{E}(g(X,Y)) = \int_\mathbb{R}\left( \int_\mathbb{R}g(x,y)d\mathbb{P}_{Y|X=x}(y) \right) d\mathbb{P}_X(x).\]
\end{theoreme}

\addcontentsline{toc}{subsection}{A noter}
\begin{remarque}[{A noter}]\phantomsection\label{III-a-noter}

\begin{itemize}
\tightlist
\item
  Ce résultat peut être interprété comme un \textbf{théorème de Fubini
  conditionnel}, dans le sens où il permet une intégration séquentielle,
  mais ici la mesure de probabilité du couple \((X,Y)\) s'exprime comme
  un produit de mesures dont l'un des termes dépend de la variable
  d'intégration de l'autre. En particulier, si on change l'ordre
  d'intégration, on change les mesures qui interviennent.
\item
  Fréquemment, dans les applications, la famille des lois
  conditionnelles est une donnée du modèle considéré, et leur existence
  ne pose donc pas de problème\,!
\item
  On retrouve les cas vus précédemment en notant que pour tout borélien
  \(B_1\) de \(\mathbb{R}\) on a
  \(\mathbb{P}_X(B_1) = \int_{B_1}d\mathbb{P}_X(x) = \sum_{x \in B_1} \mathbb{P}(X=x)\)
  lorsque \(X\) est discrète, et que pour tous boréliens \(B_1\) et
  \(B_2\) de \(\mathbb{R}\) on a
  \(\mathbb{P}_X(B_1) = \int_{B_1} f_X(x)dx\) et
  \(\mathbb{P}_{X,Y}(B_1 \times B_2) = \int_{B_1 \times B_2}f_{X,Y}(x,y) dx dy\).
\item
  Dans tout ce qui précède, les rôles de \(X\) et \(Y\) peuvent
  évidemment être inversés.
\end{itemize}
\end{remarque}

\section{Conséquences}\label{III-consuxe9quences}

Le \protect\hyperref[III-fubinicond]{théorème
précédent (p.~\pageref*{III-fubinicond})} a deux conséquences majeures.
Il fournit d'une part un moyen efficace d'identifier la loi marginale de
\(Y\) connaissant la loi marginale de \(X\) et la loi de \(Y\) sachant
\(X = x\). En effet, en notant que pour tout borélien \(B\) de
\(\mathbb{R}\),
\(\mathbb{P}_Y(B) = \mathbb{P}_{X,Y}(\mathbb{R}\times B)\) et en
appliquant ce théorème, on a la proposition suivante~:

\addcontentsline{toc}{subsection}{Formule de balayage conditionnel}
\begin{proposition}[{Formule de balayage
conditionnel}]\phantomsection\label{III-balcond}

\begin{itemize}
\tightlist
\item
  La loi marginale \(\mathbb{P}_Y\) de \(Y\) s'exprime comme la moyenne
  des lois conditionnelles \(\mathbb{P}_{Y|X=x}\) pondérée par la loi de
  \(X\). Pour tout \(B\) borélien de \(\mathbb{R}\)
  \[\mathbb{P}_Y(B) = \int_\mathbb{R}\left( \int_{B} d\mathbb{P}_{Y|X=x}(y) \right) d\mathbb{P}_X(x) = \int_\mathbb{R}\mathbb{P}_{Y|X=x}(B) d\mathbb{P}_X(x)\]
\item
  Dans le cas où \(X\) est discrète (à valeurs dans \(I\) dénombrable),
  on retrouve une expression de la formule des probabilités totales et
  composées~:
  \[\mathbb{P}_Y(B) = \mathbb{P}(Y\in B) = \sum_{x \in I} \mathbb{P}(Y \in B | X = x)\mathbb{P}(X=x)\]
\item
  Dans le cas où le couple \((X,Y)\) admet une densité, puisqu'on a
  \(f_{X,Y}(x,y) = f_{Y | X=x}(y)f_X(x)\), on obtient l'expression
  suivante pour la densité marginale~:
  \[f_Y(y) = \int_\mathbb{R}f_{X,Y}(x,y)dx = \int_\mathbb{R}f_{Y | X=x}(y)f_X(x) dx.\]
  On a en particulier la \emph{formule de Bayes pour les densités}~:
  pour tout \(x\) tel que \(f_X(x) > 0\) et tout \(y\) tel que
  \(f_Y(y) > 0\)~:
  \[ f_{X|Y=y}(x) = \frac{f_{X,Y}(x,y)}{f_Y(y)} = \frac{f_{Y|X=x}(y)f_X(x)}{f_Y(y)} .\]
\end{itemize}
\end{proposition}

\addcontentsline{toc}{subsection}{Pour fixer les idées (2)}
\begin{exemple}[{Pour fixer les idées
(2)}]\phantomsection\label{III-pour-fixer-les-iduxe9es-2}

Poursuivons \protect\hyperref[III-ex1]{l'exemple vu plus
haut (p.~\pageref*{III-ex1})}. On rappelle qu'on a déjà identifié la loi
marginale de \(X\) ainsi que la loi conditionnelle de \(Y\) sachant
\(X=n\) pour \(n\in\mathbb{N}\) que l'on rappelle ici~:
\[\mathbb{P}(X=n) = (1-\alpha)\alpha^n,\,n\in\mathbb{N}\text{ et }\forall  B \in \mathcal{B}(\mathbb{R}), \, \mathbb{P}_{Y|X=n}(B) = \int_{B \cap \mathbb{R}_+^\ast} e^{-t}\frac{t^n}{n!}dt\]
On peut en déduire la loi marginale de \(Y\) en utilisant la
\protect\hyperref[III-balcond]{formule de balayage
conditionnel (p.~\pageref*{III-balcond})} et le théorème de convergence
monotone~: \[\begin{aligned}
\mathbb{P}_Y(B)  &= \sum_{n \in \mathbb{N}} (1-\alpha)\alpha^n \int_{B \cap \mathbb{R}_+^\ast} e^{-t}\frac{t^n}{n!} dt\\
         &= (1-\alpha) \int_{B \cap \mathbb{R}_+^\ast} e^{-t} \sum_{n \in \mathbb{N}} \frac{(\alpha t)^n}{n!} dt \\
         &= \int_B 1_{\mathbb{R}_+}(t)(1-\alpha) e^{-(1-\alpha)t} dt,
\end{aligned}\] de sorte que \(Y\) suit une loi exponentielle de
paramètre \((1-\alpha)\).

En inversant les rôles, on va pouvoir identifier la loi de \(X\) sachant
\(Y \in B\) en notant que \begin{align*}
\mathbb{P}_{X,Y}(\{n\} \times B) &= \mathbb{P}_X(\{n\})\mathbb{P}_{Y|X=n}(B) \\
                         &= \int_B \frac{(\alpha t)^n}{n!} e^{-\alpha t} d\mathbb{P}_Y(t)\\
                         &= \int_B \mathbb{P}_{X | Y =t}(\{n\}) d\mathbb{P}_Y(t)
\end{align*} où l'on reconnaît que
\(\mathbb{P}_{X =n | Y =t}(\{n\}) = \frac{(\alpha t)^n}{n!} e^{-\alpha t}\),
c'est-à-dire que \(X\) sachant \(Y = t\) suit une loi de Poisson de
paramètre \(\alpha t\) pour \(\mathbb{P}_Y\)-presque tout \(t\).
\end{exemple}

En utilisant, le \protect\hyperref[III-fubinicond]{théorème de Fubini
conditionnel (p.~\pageref*{III-fubinicond})}, on obtient également une
nouvelle caractérisation de l'indépendance de deux variables aléatoires
faisant intervenir les lois conditionnelles.

\addcontentsline{toc}{subsection}{Critère d'indépendance}
\begin{proposition}[{Critère
d'indépendance}]\phantomsection\label{III-crituxe8re-dinduxe9pendance}

\begin{enumerate}
\def\labelenumi{\arabic{enumi}.}
\item
  \(X\) et \(Y\) sont indépendantes si et seulement si, pour
  \(\mathbb{P}_X\)-presque tout \(x\), \(\mathbb{P}_{Y |X = x}\) ne
  dépend pas de \(x\) et dans ce cas, on a
  \(\mathbb{P}_{Y |X = x} = \mathbb{P}_Y\), c'est-à-dire que la loi
  conditionnelle est identique à la loi marginale.
\item
  Dans le cas où \((X,Y)\) admet une densité, \(X\) et \(Y\) sont
  indépendantes si et seulement si la densité conditionnelle de \(Y\)
  sachant \(\{X = x\}\) ne dépend pas de \(x\).
\end{enumerate}
\end{proposition}

\begin{preuve}
\begin{enumerate}
\def\labelenumi{\arabic{enumi}.}
\item
  Si \(X\) et \(Y\) sont indépendantes, pour tous \(B_1\), \(B_2\)
  boréliens de \(\mathbb{R}\),
  \(\mathbb{P}_{X,Y} (B_1 \times B_2) = \mathbb{P}_X(B_1)\mathbb{P}_Y(B_2) = \int_{B_1}\mathbb{P}_Y(B_2)d\mathbb{P}_X(x) = \int_{B_2}\mathbb{P}_X(B_1) d\mathbb{P}_Y(y)\).
  Le résultat d'unicité du \protect\hyperref[III-fubinicond]{théorème de
  Fubini conditionnel (p.~\pageref*{III-fubinicond})} (à une égalité
  \(\mathbb{P}_X\)-presque sûre près), nous indique alors que
  \(\mathbb{P}_{Y|X=x}(B_2) = \mathbb{P}_Y(B_2)\).

  Inversement, si \(\mathbb{P}_{Y |X = x} = \mathbb{P}_Y\), alors
  \(\mathbb{P}_{X,Y} (B_1 \times B_2) = \int_{B_1}\mathbb{P}_{Y|X=x}(B_2)d\mathbb{P}_X(x) = \int_{B_1}\mathbb{P}_Y(B_2)d\mathbb{P}_X(x) = \mathbb{P}_X(B_1)\mathbb{P}_Y(B_2)\).
\item
  Si \(X\) et \(Y\) sont indépendantes,
  \(f_{X,Y} (x,y) = f_X(x) f_Y(y)\), d'où \(f_{Y|X=x}(y) = f_Y(y)\).

  Inversement, si \(f_{Y|X=x}(y) = f_Y(y)\) alors
  \(f_{X,Y}(x,y) = f_{Y|X=x}(y) f_X(x) = f_Y(y)f_X(x)\) et \(X\) et
  \(Y\) sont indépendantes.
\end{enumerate}

\hfill$\blacksquare$
\end{preuve}

\chapter{Espérance
conditionnelle}\label{III-espuxe9rance-conditionnelle-1}

Puisque \(\mathbb{P}_{Y|X=x}\) est la loi d'une variable aléatoire, on
peut définir l'espérance qui lui est associée et introduire la notion
d'espérance conditionnelle dans le cas où \(Y\) est intégrable.

\addcontentsline{toc}{subsection}{Espérance conditionnelle}
\begin{definition}[{Espérance
conditionnelle}]\phantomsection\label{III-defespcond}

Soit \(Y \in \mathcal{L}^1\).

\begin{enumerate}
\def\labelenumi{\arabic{enumi}.}
\tightlist
\item
  L'\emph{espérance conditionnelle de \(Y\) sachant \(\{X=x\}\)} est
  définie par
  \[\mathbb{E}(Y|X=x) = \int_\mathbb{R}y d\mathbb{P}_{Y|X=x} (y).\]
\item
  L'\emph{espérance conditionnelle de \(Y\) sachant \(X\)} est la
  \textbf{variable aléatoire} définie par~:
  \[\mathbb{E}(Y|X) = \psi(X), \text{ avec } \psi(x) = \mathbb{E}(Y|X=x).\]
\end{enumerate}
\end{definition}

\runinhead{Exercice -- Dans un triangle
(2) (\(\mathord{\bullet}\))}\phantomsection\addcontentsline{toc}{subsubsection}{Dans
un triangle (2)}\label{III-etb2}
Soient \(X\) et \(Y\) de densité jointe
\(f_{X,Y}(x,y)= \frac{1}{x}1_T (x,y)\) où \(T\) est le triangle
\(T = \{0< y< x < 1\}\). Calculer l'espérance conditionnelle de \(Y\)
sachant \(X\).\solref{III-answer-etb2}

\addcontentsline{toc}{subsection}{Conséquences}
\begin{remarque}[{Conséquences}]\phantomsection\label{III-consuxe9quences-1}

\begin{enumerate}
\def\labelenumi{\arabic{enumi}.}
\tightlist
\item
  \(\psi(x)\) n'est définie que pour \(x \notin N\), avec
  \(\mathbb{P}(X\in N) = 0\). Par conséquent, la
  \protect\hyperref[III-defespcond]{définition (p.~\pageref*{III-defespcond})}
  définit bien l'espérance conditionnelle \(\psi(X) = \mathbb{E}(Y|X)\)
  \(\mathbb{P}_X\)-presque partout, autrement dit avec probabilité 1, ou
  encore presque sûrement.
\item
  \(\mathbb{E}(\mathbb{E}(|Y||X)) = \mathbb{E}(|Y|)\) comme conséquence
  directe du \protect\hyperref[III-fubinicond]{théorème de Fubini
  conditionnel (p.~\pageref*{III-fubinicond})}. L'espérance
  conditionnelle de \(Y\) sachant \(X\) est bien définie dès que \(Y\)
  est intégrable.
\item
  Lorsque \((X,Y)\) admet une densité, l'espérance conditionnelle de
  \(Y\) sachant \(\{X=x\}\) s'écrit
  \[\mathbb{E}(Y|X=x) = \int_\mathbb{R}y f_{Y|X=x}(y) dy.\]
\end{enumerate}
\end{remarque}

\runinhead{Exercice -- Auto-conditionnement (\(\mathord{\bullet}\))}\phantomsection\addcontentsline{toc}{subsubsection}{Auto-conditionnement}\label{III-autocond}
Montrer que \(\mathbb{E}(Y|Y) = Y\).\solref{III-answer-autocond}

On peut étendre cette définition aux variables de la forme \(g(X,Y)\).

\addcontentsline{toc}{subsection}{Espérance conditionelle d'une fonction
de variables aléatoires}
\begin{definition}[{Espérance conditionelle d'une fonction de variables
aléatoires}]\phantomsection\label{III-defespcondg}

Soit \((X,Y)\) un couple de variables aléatoires réelles et \(g\) une
fonction mesurable positive ou \(\mathbb{P}_{X,Y}\)-intégrable sur
\(\mathbb{R}^2\).

\begin{enumerate}
\def\labelenumi{\arabic{enumi}.}
\tightlist
\item
  L'\emph{espérance conditionnelle de \(g(X,Y)\) sachant \(\{X=x\}\)}
  est définie par
  \[\mathbb{E}(g(X,Y)|X=x) = \int_\mathbb{R}g(x,y) d\mathbb{P}_{Y|X=x} (y).\]
\item
  L'\emph{espérance conditionnelle de \(g(X,Y)\) sachant \(X\)} est la
  \textbf{variable aléatoire} définie par~:
  \[\mathbb{E}(g(X,Y)|X) = \psi(X), \text{ avec } \psi(x) = \mathbb{E}(g(X,Y)|X=x).\]
\end{enumerate}
\end{definition}

\addcontentsline{toc}{subsection}{Espérance totale}
\begin{theoreme}[{Espérance
totale}]\phantomsection\label{III-espuxe9rance-totale}

Si \(Y\) est intégrable, alors \(\psi(X) = \mathbb{E}(Y | X)\) est
intégrable, et \[\mathbb{E}( \psi(X)) = E( Y ) .\]
\end{theoreme}

\begin{preuve}
C'est une conséquence directe du
\protect\hyperref[III-fubinicond]{théorème de Fubini
conditionnel (p.~\pageref*{III-fubinicond})}.\hfill$\blacksquare$
\end{preuve}

Ce résultat permet de calculer \(\mathbb{E}( Y )\) en conditionnant par
une variable auxiliaire \(X\)~:
\[\mathbb{E}( Y ) = \int_\mathbb{R}\mathbb{E}(Y | X = x) d\mathbb{P}_X(x)\]

Il généralise la \href{Probabilité\%20I.pdf\%20\#formprobatot}{formule
des probabilités totales}, qui correspond ici à \(Y = 1_A\) , et
\(B_x = \{X = x\}\) où les \(B_x\) forment cette fois une partition non
dénombrable de \(\mathbb{R}\). On l'écrit souvent sous la forme
\[ \mathbb{E}\left( \mathbb{E}(Y | X) \right) = \mathbb{E}( Y )\] et on
l'appelle la \emph{formule de l'espérance totale}.

L'espérance conditionnelle étant définie comme l'espérance selon la loi
conditionnelle, elle hérite des propriétés usuelles de l'espérance~:

\begin{enumerate}
\def\labelenumi{\arabic{enumi}.}
\tightlist
\item
  si \(Y\) et \(Z\) sont intégrables,
  \(\mathbb{E}(aY + bZ | X) = a \mathbb{E}(Y | X) + b \mathbb{E}(Z | X)\),
\item
  \(\mathbb{E}(Y | X) \geq 0\) si \(Y \geq 0\),
\item
  \(\mathbb{E}(1 | X) = 1\).
\end{enumerate}

De plus, si \(g\) est mesurable positive ou \(\mathbb{P}_X\)-intégrable,
\[ \mathbb{E}(Y g(X) | X) = g(X) \mathbb{E}(Y | X) \] est une
généralisation de l'égalité 1. ci-dessus, au cas où \(a = g(X)\), qui
doit être considéré ``comme une constante'' dans le calcul de
l'espérance conditionnelle sachant \(X\) (\(X\) est fixée comme une
donnée connue a priori). En effet, on a alors
\(\mathbb{E}(g(x)Y|X=x) = g(x)\psi(x)\). Enfin, on déduit directement du
\protect\hyperref[III-fubinicond]{théorème de Fubini
conditionnel (p.~\pageref*{III-fubinicond})} la proposition suivante.

\addcontentsline{toc}{subsection}{Transfert conditionnel}
\begin{proposition}[{Transfert
conditionnel}]\phantomsection\label{III-transfcond}

Soient un couple \((X,Y)\) de variables aléatoires réelles de loi jointe
\(\mathbb{P}_{X,Y}\) et \(g\) une fonction mesurable positive ou
\(\mathbb{P}_{X,Y}\)-intégrable sur \(\mathbb{R}^2\). On a pour
\(\mathbb{P}_X\)-presque tout \(x\) dans \(\mathbb{R}\)
\[\mathbb{E}(g(X,Y)|X=x) = \mathbb{E}(g(x,Y)|X=x) = \int_{\mathbb{R}}g(x,y) d\mathbb{P}_{Y|X=x} (y)\]
Si de plus \(X\) et \(Y\) sont indépendantes, on a~:
\[\mathbb{E}(g(X,Y)|X=x) = \mathbb{E}(g(x,Y)|X=x) = \int_{\mathbb{R}}g(x,y) d\mathbb{P}_Y(y).\]

Autrement dit, lorsqu'on conditionne par l'événement \(\{X=x\}\), cela
revient à fixer la valeur de la variable aléatoire \(X\) à la constante
\(x\).
\end{proposition}

\runinhead{Exercice -- Espérance conditionnelle d'un produit de
variables (\(\mathord{\bullet}\))}\phantomsection\addcontentsline{toc}{subsubsection}{Espérance
conditionnelle d'un produit de variables}\label{III-prod}
Calculer \(\mathbb{E}(XY|X=x)\) puis
\(\mathbb{E}(XY|X)\).\solref{III-answer-prod}

\addcontentsline{toc}{subsection}{Vecteurs gaussiens à densité}
\subsection*{Exemple~: vecteurs gaussiens à
densité}\label{III-exemple-vecteurs-gaussiens-uxe0-densituxe9}

Dans ce qui précède, on a décrit les lois et les espérances
conditionnelles dans le cas d'un couple de variables aléatoires à
valeurs dans \(\mathbb{R}^2\). Ces résultats sont aussi valables pour
des couples de vecteurs, dont on décrit ici un cas particulier.

Dans le cas des vecteurs gaussiens à densité, c'est-à-dire dont la
matrice de covariance est définie positive et donc inversible, le calcul
des lois conditionnelles de certaines composantes par rapport aux autres
est particulièrement aisé. On va voir en particulier que les lois
conditionnelles ont le bon goût d'être elles-mêmes gaussiennes, ce qui
explique (en partie) le succès de ces modèles dans les applications.

On considère un vecteur gaussien \(X = (X_1,\ldots,X_n)\) à valeurs dans
\(\mathbb{R}^n\) d'espérance \(m\) et de matrice de covariance \(C\)
définie positive. On a vu au chapitre 2 que la densité du vecteur \(X\)
s'écrit pour \(x\in\mathbb{R}^d\)~:
\[f_X(x) = \frac{1}{(2\pi)^{n/2}\sqrt{\det (C)}}\exp \left(-\frac{1}{2}(x-m)^t C^{-1}(x-m)\right)\]

Soit \(1 \leq k < n\) un entier. On souhaite exprimer \(f_{Y|Z=z}\), la
densité conditionnelle de \(Y = (X_1,\ldots,X_k)\) sachant
\(Z = (X_{k+1},\ldots,X_n) = (x_k,\ldots,x_n) = z\) (si \(k+1 = n\), ce
vecteur se réduit à une seul valeur). On a vu que
\[f_X = f_{Y|Z=z} f_Z,\] où \(f_Z\) est la densité marginale de \(Z\).
On cherche donc à décomposer \(f_X\) de la sorte. On note
\(m = (m_Y,m_Z)\) et on remarque que \(C\) peut se décomposer en blocs~:
\begin{equation*}
C = \left(\begin{array}{cc} C_Y & C_{Y,Z} \\ C_{Z,Y} & C_Z \end{array}\right)
\end{equation*} où \(C_Y = \text{Cov}(Y,Y)\), \(C_Z = \text{Cov}(Z,Z)\)
et \(C_{Y,Z} = \text{Cov}(Y,Z)\). Le \emph{complément de
Schur}\footnote{voir par exemple l'excellent
  \href{https://www2.compute.dtu.dk/pubdb/edoc/imm3274.pdf}{matrix
  cookbook}.} du bloc \(C_Y\) est la matrice
\[CS_Y = C_Y - C_{Y,Z}C_Z^{-1}C_{Z,Y}\] et permet d'exprimer l'inverse
de \(C\) comme~: \begin{equation*}
C^{-1} = \left(\begin{array}{cc} CS_Y^{-1} & -CS_Y^{-1}C_{Y,Z}C_Z^{-1} \\ -C_Z^{-1}C_{Z,Y}CS_Y^{-1}  & C_Z^{-1} + C_Z^{-1}C_{Z,Y}CS_Y^{-1}C_{Y,Z}C_Z^{-1} \end{array}\right)
\end{equation*} On peut alors réarranger les termes de la forme
quadratique dans \(f_X\) et on obtient~: \begin{align*}
(x-m)^t C^{-1}(x-m) =& \left(y - (m_Y + C_{Y,Z}C_Z^{-1}(z-m_Z))\right)^t CS_Y^{-1}\\
&.\left(y - (m_Y + C_{Y,Z}C_Z^{-1}(z-m_Z))\right)\\
 &+ (z-m_Z)^t C_Z^{-1}(z-m_Z)
\end{align*} Pour la constante, on peut remarquer que~:
\[ \det (C) = \det(CS_Y)\det(C_Z).\] On en déduit ainsi que
\[f_{Y|Z=z}(y) = \frac{1}{(2\pi)^{k/2}\sqrt{\det (CS_Y)}}\exp \left(-\frac{1}{2}\left(y - \psi(z)\right)^t CS_Y^{-1}\left(y - \psi(z))\right)\right)\]

C'est-à-dire que la variable aléatoire \(Y|Z=z\) est gaussienne
d'espérance \(m_{Y|Z=z} = \psi(z) = m_Y + C_{Y,Z}C_Z^{-1}(z-m_Z)\) et de
matrice de covariance \(CS_Y = C_Y - C_{Y,Z}C_Z^{-1}C_{Z,Y}\). Autrement
dit, l'espérance conditionnelle de \(Y\) sachant \(Z\) est la variable
aléatoire \(\mathbb{E}(Y|Z) = \psi(Z) =(m_Y + C_{Y,Z}C_Z^{-1}(Z-m_Z))\).
On notera que la covariance conditionnelle donnée par \(CS_Y\) ne dépend
pas de la valeur prise par \(Z\).

\chapter{Régression et espérance conditionnelle des variables de carré
intégrable}\label{III-ruxe9gression-et-espuxe9rance-conditionnelle-des-variables-de-carruxe9-intuxe9grable}

La régression est un ensemble de méthodes (d'apprentissage) statistiques
très utilisées pour analyser la relation d'une variable par rapport à
une ou plusieurs autres. Ces méthodes visent notamment à décrire les
liens de dépendance entre variables mais aussi à prédire au mieux la
valeur d'une quantité non observée en fonction d'une ou plusieurs autres
variables.

Pour fixer les idées, supposons que l'on s'intéresse au lien entre une
variable d'intérêt \(Y\) et une variable explicative \(X\) dans le but
de prédire la valeur de \(Y\) étant donnée l'observation de \(X\). Le
but est alors de proposer une certaine fonction
\(f : \mathbb{R}\rightarrow \mathbb{R}\) qui permette d'approximer \(Y\)
par \(f(X)\). Une telle fonction \(f\) est appelée un \emph{prédicteur}
de \(Y\) en fonction de \(X\).

On va montrer dans cette section que, dans le cadre des variables de
carré intégrable, une solution naturelle à ce problème est précisément
donnée par l'espérance conditionnelle de \(Y\) sachant \(X\). Nous nous
appuierons pour ce faire sur une interprétation géométrique de
l'espérance conditionnelle en tant que projection orthogonale dans un
espace de Hilbert. Cette approche justifie en particulier l'usage des
méthodes de moindres carrés dans les problèmes de régression.

\section{Rappels sur les espaces de Hilbert et application en
probabilités}\label{III-rappels-sur-les-espaces-de-hilbert-et-application-en-probabilituxe9s}

Soit \((\Omega, \mathcal{A},\mathbb{P})\) un espace probabilisé. On
rappelle que par définition, une variable aléatoire \(Y\) sur cet espace
est une fonction \(\mathcal{A}\)-mesurable
\(Y : \Omega \rightarrow \mathbb{R}\). On considère alors l'espace
fonctionnel \(L^2(\Omega)\) introduit dans le cours de calcul intégral.
Pour rappel, cet espace contient l'ensemble des fonctions
\(\mathcal{A}\)-mesurables (i.e.~des variables aléatoires) \(Y\) telles
que \(\vert Y\vert^2\) soit intégrable c'est-à-dire telles que \[
\Vert Y\Vert_2 = \bigg(\int_{\Omega} \vert Y(\omega)\vert^2 d\mathbb{P}(\omega)\bigg)^{1/2}=\mathbb{E}(Y^2)^{1/2} < \infty.
\] En particulier, contrairement à l'espace \(\mathcal{L}^2\) introduit
au chapitre précédent, dans l'espace \(L^2(\Omega)\) les fonctions
égales \(\mathbb{P}\)-presque partout sont confondues. C'est-à-dire que
deux fonctions \(X\) et \(Y\) égales \(\mathbb{P}\)-presque partout
correspondent à un même élément de \(L^2(\Omega)\) appelé classe
d'équivalence. Ainsi, pour \(X, Y \in L^2(\Omega)\), l'égalité \(X=Y\)
signifiera donc que \(X\) et \(Y\) sont dans la même classe
d'équivalence, autrement dit que les variables \(X\) et \(Y\) sont
égales presque sûrement.

On a vu dans le cours de calcul intégral que l'espace \(L^2(\Omega)\),
muni de la norme \(\Vert \cdot\Vert_2\), était un espace de Banach,
c'est-à-dire un espace vectoriel normé et complet pour cette norme. On
va maintenant introduire un produit scalaire sur cet espace, dont la
norme associée sera précisément la norme \(\Vert \cdot\Vert_2\), et
montrer que \(L^2(\Omega)\) est ainsi un espace de Hilbert.

On considère à cet effet l'application bilinéaire
\(\langle \cdot \vert \cdot\rangle\) définie par
\[\langle X \vert Y\rangle = \int_{\Omega} X(\omega) Y(\omega) d\mathbb{P}(\omega) = \mathbb{E}(XY), \quad X,Y\in L^2(\Omega).\]
L'application \(\langle \cdot \vert \cdot\rangle\) est bien définie
puisque, d'après l'inégalité de Cauchy-Schwartz, on a pour tous
\(X,Y\in L^2(\Omega)\):
\[\langle X \vert Y\rangle^2=\mathbb{E}(XY)^2 \leq \mathbb{E}(X^2)\mathbb{E}(Y^2)=\Vert X\Vert_2^2\Vert Y\Vert_2^2<\infty.\]
De plus, l'application \(\langle \cdot \vert \cdot\rangle\) définit un
produit scalaire sur \(L^2(\Omega)\) car elle vérifie les propriétés
suivantes pour tous \(X,Y\in L^2(\Omega)\):

\begin{itemize}
\tightlist
\item
  Symétrique:
  \(\langle X \vert Y\rangle= \mathbb{E}(XY)=\mathbb{E}(YX)=\langle X \vert Y\rangle\),
\item
  Positive: \(\langle X \vert X\rangle= \mathbb{E}(X^2) \ge 0\),
\item
  Définie:
  \(\langle X \vert X\rangle=0 \Leftrightarrow \mathbb{E}(X^2)=0 \Leftrightarrow X=0\)
  .
\end{itemize}

On remarque enfin que la norme associée au produit scalaire
\(\langle \cdot \vert \cdot\rangle\) est bien la norme
\(\Vert\cdot\Vert_2\) introduite précédemment (car pour tout
\(X\in L^2(\Omega)\),
\(\langle X \vert X\rangle^{1/2}=\mathbb{E}(X^2)^{1/2}=\Vert X\Vert_2\)).
On en déduit que l'espace \(L^2(\Omega)\) muni du produit scalaire
\(\langle \cdot \vert \cdot\rangle\) est un espace de Hilbert.

L'intérêt de travailler avec un espace de Hilbert est qu'il est possible
d'y définir de manière univoque la notion de projection orthogonale. En
effet, la projection orthogonale d'un élément \(Y\in L^2(\Omega)\) sur
un sous-espace convexe, fermé et non-vide \(K \subset L^2(\Omega)\) est
définie comme l'unique élément \(P_K(Y)\in K\) qui minimise la distance
\(\Vert Y - P_K(Y)\Vert_2\), autrement dit l'erreur des moindres carrés
\(\mathbb{E}((Y-P_K(Y))^2)\) entre \(Y\) et les éléments de \(K\). On
retrouve en particulier l'objectif de prédiction décrit en introduction,
ce qui permet de conclure que la projection orthogonale \(P_K(Y)\) est
le prédicteur optimal de \(Y\) au sens des moindres carrés parmi les
éléments de \(K\). On en déduit la propriété suivante qui permet de
caractériser ce prédicteur optimal d'après les propriétés de la
projection orthogonale vues dans le cours de calcul intégral.

\addcontentsline{toc}{subsection}{Prédicteur optimal sur un convexe
fermé}
\begin{proposition}[{Prédicteur optimal sur un convexe
fermé}]\phantomsection\label{III-identEC}

Soit \(Y\in L^2(\Omega)\) et soit \(K \subset L^2(\Omega)\) convexe,
fermé et non-vide. Le prédicteur optimal, au sens des moindres carrés,
de \(Y\) sur \(K\) est la projection orthogonale de \(Y\) sur \(K\). Il
est donc caractérisé comme étant l'unique élément de \(K\) vérifiant la
propriété
\[ \langle Y - P_K(Y) \vert U - P_K(Y)\rangle \le 0 \quad  \forall ~U\in K.\]
Si de plus \(K\) est un sous-espace vectoriel, cette caractérisation
prend une forme plus simple, à savoir:
\[ \langle Y - P_K(Y) \vert U \rangle = 0 \quad  \forall ~U\in K.\]
\end{proposition}

Pour revenir à notre problème initial de prédiction optimale de \(Y\)
par une fonction de \(X\), il s'agit donc de choisir le sous-espace
\(K\) de projection de manière à ce qu'il ne contienne que des éléments
s'exprimant comme des fonctions de \(X\), et d'identifier parmi les
éléments de \(K\) celui qui satisfait la
\protect\hyperref[III-identEC]{propriété
ci-dessus (p.~\pageref*{III-identEC})}.

\section{Cas particulier: la régression
linéaire}\label{III-cas-particulier-la-ruxe9gression-linuxe9aire}

Reprenons les variables \(X,Y\in L^2(\Omega)\) introduites en début de
section, et le problème initial de prédiction de \(Y\) à partir de
\(X\). Dans cette sous-section, nous nous restreignons à des prédicteurs
affines, c'est-à-dire que nous souhaitons trouver la meilleure
approximation, au sens des moindres carrés, de \(Y\) par une fonction
affine de \(X\) de la forme \(f(X)=aX + b\). Un tel prédicteur est
appelé \emph{régression linéaire de \(Y\) par \(X\)}.

Tout prédicteur affine \(f(X)=aX+b\) peut être vu comme une combinaison
linéaire de deux variables aléatoires de \(L^2(\Omega)\)~: \(X\), et la
variable aléatoire \(1\) égale à \(1\) presque sûrement. Par conséquent,
\(f(X)\in K\) où \(K\) désigne le sous-espace vectoriel de
\(L^2(\Omega)\) engendré par \(X\in L^2(\Omega)\) et
\(1\in L^2(\Omega)\). Ainsi, trouver un prédicteur affine optimal de
\(Y\) revient à trouver un prédicteur optimal de \(Y\) dans le
sous-espace vectoriel \(K\).

\(K\) étant un sous-espace vectoriel fermé de \(L^2(\Omega)\) (car de
dimension finie), la \protect\hyperref[III-identEC]{propriété de
caractérisation (p.~\pageref*{III-identEC})} du prédicteur optimal
permet de conclure que le prédicteur affine optimal de \(Y\) par \(X\)
est la projection orthogonale \(P_K(Y)\) de \(Y\) sur \(K\). En
particulier, puisque tout élément de \(K\) est une combinaison linéaire
de \(X\) et \(1\), on peut écrire que \(P_K(Y) = a^*X + b^* 1\), où les
coefficients \(a^*,b^*\) vérifient :
\[ \langle Y - (a^*  X + b^*  1) \vert a  X + b  1 \rangle = 0 \quad  \forall a,b\in\mathbb{R},\]
Cette égalité devant être vérifiée pour tous \(a,b\in\mathbb{R}\), elle
implique, par bilinéarité du produit scalaire, que les coefficients
\(a^*\) et \(b^*\) soient solutions du système : \[
\left\lbrace\begin{aligned}
   &\langle Y \vert X \rangle - a^*\langle X \vert X \rangle-b^*\langle 1 \vert X \rangle  = 0 \\
   &\langle Y \vert 1 \rangle - a^*\langle X \vert 1 \rangle-b^*\langle 1 \vert 1 \rangle = 0
\end{aligned}\right.
\] Par définition du produit scalaire \(\langle\cdot\vert\cdot\rangle\),
ce système se réécrit sous la forme : \[
\left\lbrace\begin{aligned}
   &\mathbb{E}(XY) = a^*\mathbb{E}(X^2)+b^*\mathbb{E}(X)   \\
   &\mathbb{E}(Y) = a^*\mathbb{E}(X)+b^* 
\end{aligned}\right.
\] et a pour solution \[\begin{aligned}
a^* & = \frac{\mathbb{E}(XY)-\mathbb{E}(X)\mathbb{E}(Y)}{\mathbb{E}(X^2)-\mathbb{E}(X)^2}=\frac{\text{Cov}(X,Y)}{\mathbb{V}(X)} = \rho(X,Y)\frac{\sigma_Y}{\sigma_X} \\
b^* & = \mathbb{E}(Y)  - a^*~\mathbb{E}(X)
\end{aligned}\]

En conclusion, le meilleur prédicteur affine, au sens des moindres
carrés, de \(Y\) par \(X\) est la fonction \(f^*(X)\) définie par
\[ f^*(X)=P_K(Y)=a^* X + b^*=\mathbb{E}(Y) + \rho(X,Y)\frac{\sigma_Y}{\sigma_X} (X -\mathbb{E}(X)),\]
et l'erreur quadratique moyenne associée à ce prédicteur vaut :
\[\begin{aligned}
\mathbb{E}\left(\left(Y - f^*(X)\right)^2\right) & = \sigma_Y^2 + \rho^2(X,Y)\sigma^2_Y - 2\rho^2(X,Y)\sigma^2_Y\\
                        & = \sigma^2_Y(1-\rho^2(X,Y)).
\end{aligned}\] On voit ainsi que cette erreur est proche de 0 lorsque
\(|\rho(X,Y)| \approx 1\) tandis qu'elle est proche de
\(\mathbb{V}(Y) = \sigma^2_Y\) lorsque \(\rho(X,Y) \approx 0\). On
notera au passage qu'on obtient que la meilleure approximation de \(Y\)
par une constante (ce qui revient à prendre \(X=0\)) est son espérance.

\runinhead{Exercice -- Une deuxième
approche (\(\mathord{\bullet}\))}\phantomsection\addcontentsline{toc}{subsubsection}{Une
deuxième approche}\label{III-detail}
Montrer que les coefficients \(a^*\) et \(b^*\) sont les uniques
minimiseurs de l'erreur \(\ell(a,b)=\mathbb{E}((Y - (aX + b))^2)\), pour
\(a,b\in\mathbb{R}\).\solref{III-answer-detail}

\addcontentsline{toc}{subsection}{Limites de la régression linéaire}
\begin{remarque}[{Limites de la régression
linéaire}]\phantomsection\label{III-limites-de-la-ruxe9gression-linuxe9aire}

L'hypothèse d'une relation linéaire est très forte et pas nécessairement
toujours adaptée pour expliquer des relations de dépendances entre
variables. Soit en effet une variable aléatoire réelle \(X\) de
\(\mathcal{L}^3\) (i.e.~\(X^3\) est \(\mathbb{P}_X\) intégrable)
symétrique, c'est-à-dire telle que \(X\) et \(-X\) sont de même loi. On
a alors \(\mathbb{E}(X) = -\mathbb{E}(X) = 0\). Les variables \(X\) et
\(X^2\) ne sont clairement pas indépendantes. Pour autant, on a
\(\text{Cov}(X,X^2) = \mathbb{E}(X^3) = -\mathbb{E}(X^3) = 0\) et le
coefficient de régression \(a^*\) est nul.
\end{remarque}

\section{Cas général}\label{III-cas-guxe9nuxe9ral-1}

Dans le paragraphe précédent, on s'est intéressé à approximer une
variable aléatoire \(Y\in L^2(\Omega)\) par une fonction affine d'une
autre variable \(X\in L^2(\Omega)\). On va montrer ici que si on
généralise la recherche du prédicteur optimal à l'ensemble des fonctions
\(f\) telles que \(f(X)\) est de carré intégrable, la solution est
donnée par l'espérance conditionnelle \(\psi(X) = \mathbb{E}(Y|X)\).

En effet, considérons le sous-espace
\(L^2_X(\Omega)\subset L^2(\Omega)\) constitué des (classes
d'équivalence) des variables aléatoires de la forme \(f(X)\), avec
\(f : \mathbb{R}\rightarrow\mathbb{R}\) borélienne et telle que
\(f(X) \in L^2(\Omega)\). On remarque que \(L^2_X(\Omega)\) est
sous-espace vectoriel de \(L^2(\Omega)\). On admettra par ailleurs que
\(L^2_X(\Omega)\) est fermé. Par conséquent, le prédicteur optimal de
\(Y\) dans \(L^2_X(\Omega)\) est donné par la projection orthogonale de
\(Y\) sur \(K=L^2_X(\Omega)\), et on peut une nouvelle fois appliquer la
\protect\hyperref[III-identEC]{propriété de
caractérisation (p.~\pageref*{III-identEC})} associée pour l'identifier.

On note par ailleurs que l'espérance conditionnelle
\(\psi(X) = \mathbb{E}(Y|X) \in L^2_X(\Omega)\) et que pour tout
\(f(X)\in L^2_X(\Omega)\), \[ 
\begin{aligned}
  \langle Y - \psi(X) \vert f(X) \rangle =  \mathbb{E}\big((Y - \psi(X))f(X)\big)
&= \mathbb{E}\big(Yf(X)\big) - \mathbb{E}\big(\mathbb{E}(Y|X)f(X)\big) \\
&= \mathbb{E}\big(Yf(X)\big) - \mathbb{E}\big(\mathbb{E}(Yf(X)|X)\big) \\
&= \mathbb{E}\big(Yf(X)\big) - \mathbb{E}\big(Yf(X)\big)=0
\end{aligned}
\] où la seconde égalité se déduit du fait que \(f(X)\) est
\(\mathbb{P}\)-intégrable, et la troisième découle du théorème de
l'espérance totale. Ainsi l'espérance conditionnelle vérifie la
\protect\hyperref[III-identEC]{propriété de
caractérisation (p.~\pageref*{III-identEC})} et par unicité de la
projection orthogonale, on en déduit que
\(P_{L^2_X(\Omega)}(Y)=\psi(X)\).

\begin{figure}
\centering
\includegraphics{Proj.tex.pdf}
\caption{Illustration géométrique de l'espérance conditionnelle comme
projection orthogonale sur un sous-espace de
\(L^2(\Omega)\)}\label{III-fig_proj}
\end{figure}

En conclusion, l'espérance conditionnelle \(\mathbb{E}(Y|X)\) est le
meilleur prédicteur de \(Y\) au sens des moindres carrés parmi les
fonctions de \(X\). On notera qu'en pratique, dans les problèmes de
régression, on se restreint souvent à des classes de fonctions
paramétrées (par exemple les fonctions affines vues précédemment, ou les
réseaux de neurones dans l'apprentissage profond) car on ne dispose que
rarement de l'expression explicite de l'espérance conditionnelle.

Il est ensuite immédiat que le ``résidu'' \(Y-\mathbb{E}(Y|X)\) est non
corrélé avec \(X\) du fait de l'orthogonalité. On en déduit ainsi la
\emph{formule de la variance totale}~: \[
\begin{aligned}
\mathbb{V}(Y) = \|Y - \mathbb{E}(Y)\|^2 &=  \|Y - \mathbb{E}(Y|X) + \mathbb{E}(Y|X) - \mathbb{E}(Y)\|^2 \\
                          &=  \|Y - \mathbb{E}(Y|X)\|^2 + \|\mathbb{E}(Y|X) - \mathbb{E}(Y)\|^2 \\
                          &=  \mathbb{E}((Y - \mathbb{E}(Y|X))^2) + \mathbb{E}((\mathbb{E}(Y|X) - \mathbb{E}(Y))^2)\\
                          &= \mathbb{E}(\mathbb{E}((Y - \mathbb{E}(Y|X))^2|X)) + \mathbb{V}(\mathbb{E}(Y|X))\\
                          &= \mathbb{E}(\mathbb{V}(Y|X)) + \mathbb{V}(\mathbb{E}(Y|X)).
\end{aligned}
\] où on a utilisé la formule de l'espérance totale et introduit la
variable aléatoire variance conditionnelle
\(\mathbb{V}(Y|X) = \mathbb{E}((Y - \mathbb{E}(Y|X))^2|X)\) comme cas
particulier de la \protect\hyperref[III-defespcondg]{définition vue plus
haut (p.~\pageref*{III-defespcondg})}. Cette décomposition est illustrée
dans la \protect\hyperref[III-fig_proj2]{figure
ci-après (figure~\ref*{III-fig_proj2})}.

\begin{figure}
\centering
\includegraphics{Proj2.tex.pdf}
\caption{Illustration géométrique de la formule de la variance totale.
L'espace \(\mathrm{Vect}\lbrace 1\rbrace\) désigne le sous-espace de
\(L^2(\Omega)\) composé des variables aléatoires p.s.
constantes.}\label{III-fig_proj2}
\end{figure}

\runinhead{Exercice -- Variance
totale (\(\mathord{\bullet}\))}\phantomsection\addcontentsline{toc}{subsubsection}{Variance
totale}\label{III-vartot}
Démontrer ce résultat sans utiliser la notion
d'orthogonalité.\solref{III-answer-vartot}

\chapter{Exercices}\label{III-exercices}

\section{Couple de variables}\label{III-couple-de-variables}

Soient \(X\) et \(Y\) deux v.a. réelles. On suppose que la densité
conditionnelle de \(X\) sachant \(Y = y\) est la densité
\(1_{\mathbb{R}_+}(x)y^2 xe^{-xy}\) et que la loi de \(Y\) est de
densité \(\frac{1}{y^2} 1_{[1,+\infty[} (y)\). On pose \(T = XY\).

\runinhead{Question
1}\phantomsection\addcontentsline{toc}{subsubsection}{Question
1}\label{III-ex1-1}
Trouver la loi du couple \((T,Y)\). Qu'en
déduit-on\,?\solref{III-answer-ex1-1}

\runinhead{Question
2}\phantomsection\addcontentsline{toc}{subsubsection}{Question
2}\label{III-ex1-2}
Trouver la loi conditionnelle de \(Y\) sachant
\(X = x\).\solref{III-answer-ex1-2}

\runinhead{Question
3}\phantomsection\addcontentsline{toc}{subsubsection}{Question
3}\label{III-ex1-3}
Calculer \(E(Y |X)\).\solref{III-answer-ex1-3}

\section{Mélanges de lois}\label{III-muxe9langes-de-lois}

\emph{Adapté du cours de probabilités de S. Bonnabel et M. Schmidt
(MINES ParisTech).}

Pour modéliser un phénomène multimodal, on utilise souvent des mélanges
de gaussiennes. C'est le cas notamment en classification non-supervisée,
où on fait l'hypothèse que chacune des classes suit une loi gaussienne.
Soient \(n\in\mathbb{N}^\ast\) et \(K\) une variable aléatoire prenant
les valeurs \(1,\dots,n\) avec les probabilités non nulles
\(p_1,\dots,p_n\) telles que \(\sum_{i = 1}^n p_i = 1\). Soient
\(X_1,\dots,X_n\) des variables aléatoires gaussiennes mutuellement
indépendantes, d'espérances respectives \(m_1,\dots,m_n \in \mathbb{R}\)
et de variances respectives
\(\sigma_1^2,\dots,\sigma_n^2 \in\mathbb{R}_+^\ast\), toutes
indépendantes de \(K\). On appelle mélange de gaussiennes la loi de la
variable aléatoire \(X = X_K\). Pour tout \(i\in\{1,\dots,n\}\), on
notera \(f_i\) la densité de la variable aléatoire \(X_i\).

\runinhead{Question
1}\phantomsection\addcontentsline{toc}{subsubsection}{Question
1}\label{III-melloi1}
Soit \(i \in \{1,\dots,n\}\). Quelle est la densité \(f_{X\mid K = i}\)
de \(X\) conditionnellement à l'événement
\(\{K = i\}\)\,?\solref{III-answer-melloi1}

\runinhead{Question
2}\phantomsection\addcontentsline{toc}{subsubsection}{Question
2}\label{III-melloi2}
Calculer la densité de probabilité de la variable
\(X\).\solref{III-answer-melloi2}

\runinhead{Question
3}\phantomsection\addcontentsline{toc}{subsubsection}{Question
3}\label{III-melloi3}
Calculer \(\mathbb{E}(X)\). Montrer que
\(\mathbb{V}(X) = \sum_{i = 1}^n p_i\sigma_i^2 + \bar{\sigma}^2\), où ce
dernier terme peut être interprété comme la dispersion des
espérances.\solref{III-answer-melloi3}

\runinhead{Question
4}\phantomsection\addcontentsline{toc}{subsubsection}{Question
4}\label{III-melloi4}
Comment approximeriez-vous le mélange par une unique gaussienne\,? Faire
un schéma dans le cas \(m = 2\).\solref{III-answer-melloi4}

\section{Lois conjuguées}\label{III-lois-conjuguuxe9es}

Soit un vecteur aléatoire \((X,Y)\) de loi jointe \(\mathbb{P}_{X,Y}\).
Expliciter la loi conditionnelle de \(Y\) sachant \(\{X=x\}\) dans les
situations suivantes, en prenant soin d'expliciter pour quelles valeurs
de \(x\) ces dernières ont du sens.

\runinhead{Question
1}\phantomsection\addcontentsline{toc}{subsubsection}{Question
1}\label{III-loiconj-expexp}
\(Y\) suit une loi Exponentielle de paramètre
\(\lambda \in\mathbb{R}_+^\ast\) et pour tout \(y\in\mathbb{R}_+^\ast\),
la variable aléatoire \(X\) sachant \(\{Y=y\}\) suit une loi
Exponentielle de paramètre \(y\).\solref{III-answer-loiconj-expexp}

\runinhead{Question
2}\phantomsection\addcontentsline{toc}{subsubsection}{Question
2}\label{III-loiconj-gampoi}
\(Y\) suit une loi Gamma de paramètres
\(\alpha,\theta \in \mathbb{R}_+^\ast\) et pour tout
\(y\in\mathbb{R}_+^\ast\), la variable aléatoire \(X\) sachant
\(\{Y=y\}\) suit une loi de Poisson de paramètre
\(y\).\solref{III-answer-loiconj-gampoi}

\section{Randomisation}\label{III-randomize}

\emph{Extrait du cours de probabilités de S. Bonnabel et M. Schmidt
(MINES ParisTech).}

Des clients arrivent à la boutique SNCF du boulevard Saint-Michel à des
instants aléatoires. On note \(T_0\) l'heure d'ouverture puis
\(T_1, T_2, \dots\) les temps successifs d'arrivée des clients jusqu'à
l'heure de fermeture. Les études statistiques montrent qu'on peut, dans
une tranche horaire donnée, supposer que les temps d'attente
\(X_1 = T_1 -T_0, X_2 = T_2 -T_1,\dots\) peuvent être modélisés par des
variables aléatoires indépendantes et de même loi qu'une variable
aléatoire positive \(X\). Par ailleurs, une loterie interne décide que
chaque jour dans la tranche horaire considérée, le \(N^{\text{ème}}\)
client sera l'heureux gagnant d'un trajet gratuit Paris-La Ciotat, où
\(N\) est une variable aléatoire bornée dont la loi dépend du processus
de loterie (e.g.~tous les clients entre le premier et le
\(30^{\text{ème}}\) ont une chance \(1/30\) d'être tirés au sort, en
supposant qu'on est sûr d'avoir au moins \(30\) clients dans la tranche
horaire).

On se demande alors~: quel est le temps d'attente moyen avant d'obtenir
un gagnant\,?\solref{III-answer-randomize}

\section{Etats cachés --- indépendance
conditionnelle}\label{III-etats-cachuxe9s-induxe9pendance-conditionnelle}

Soucieux de l'évolution du potager de l'école, des élèves à la main
verte s'intéressent à l'évolution de la température dans le jardin côté
Luxembourg. Ils récupèrent pour cela un thermomètre dans un laboratoire,
l'installent près du potager, et en relèvent les mesures à intervalles
de temps réguliers. Les résultats les surprennent rapidement~: les
températures affichées ne correspondent pas à celles prévues par
météo-France. Leur thermomètre est sans doute déréglé.

On se propose de les aider à comprendre le phénomène dont ils sont
témoins à l'aide d'un modèle probabiliste particulier, nommé
\emph{modèle de Markov caché}. Précisément, on considère la suite des
vraies températures que l'on aurait souhaité relever comme une suite de
v.a.r. non indépendantes \((X_n)_{n\in\mathbb{N}^\ast}\), dite
\emph{d'états cachés} (on ne les observe pas directement). Les erreurs
commises par le thermomètre sont quant à elles modélisées par une suite
de v.a.r. \((\epsilon_n)_{n\in\mathbb{N}^\ast}\), toutes indépendantes
et de même loi admettant une densité \(f_\epsilon\). Elles sont
supposées indépendantes de la suite \((X_n)_{n\in\mathbb{N}^\ast}\)
(l'erreur du thermomètre lui est propre et ne dépend pas de la
température réelle). A chaque instant \(n\in\mathbb{N}^\ast\), on
suppose que la mesure du thermomètre est la variable aléatoire
\[Y_n = X_n + \epsilon_n,\] et que le vecteur aléatoire
\((X_1,\dots,X_n)\) possède une densité jointe notée \(f_{1:n}\).

\runinhead{Question
1}\phantomsection\addcontentsline{toc}{subsubsection}{Question
1}\label{III-ec1}
Montrer que pour tout \(n \in \mathbb{N}^\ast\) et tout
\(x\in\mathbb{R}\), la loi de \(Y_n\) sachant \(\{X_n = x\}\) admet une
densité, que l'on explicitera.\solref{III-answer-ec1}

\runinhead{Question
2}\phantomsection\addcontentsline{toc}{subsubsection}{Question
2}\label{III-ec2}
Montrer que les \(n\in\mathbb{N}^\ast\) relevés de température
\(Y_1,\dots,Y_n\) sont indépendants \textbf{conditionnellement} aux
états cachés \(X_1,\dots,X_n\).\solref{III-answer-ec2}

\section{Covariance totale}\label{III-covtot}

Soient \(X\), \(Y\) et \(Z\) trois variables aléatoires réelles de carré
intégrable. La covariance conditionnelle de \(X\) et \(Y\) sachant \(Z\)
est définie comme la variable aléatoire
\[\text{Cov}(X,Y \mid Z) = \mathbb{E}\Bigl( \bigl( X - \mathbb{E}(X\mid Z) \bigr)\bigl( Y - \mathbb{E}(Y\mid Z) \bigr) \Bigm| Z  \Bigr).\]

Etablir la formule de la covariance totale~:
\[\text{Cov}(X,Y) = \mathbb{E}\bigl(\text{Cov}(X,Y\mid Z)\bigr) + \text{Cov}\bigl( \mathbb{E}(X\mid Z), \mathbb{E}(Y\mid Z) \bigr).\]

\solref{III-answer-covtot}

\section{Débruitage (formule de
Tweedie)}\label{III-duxe9bruitage-formule-de-tweedie}

Les photographies sont souvent entâchées d'un bruit qui dégrade leur
qualité, notamment lorsque la lumière est faible. Pour améliorer le
rendu, des algorithmes de débruitage sont mis en oeuvre. Dans cet
exercice, on se propose de donner un éclairage probabiliste sur ce
problème.

On considère que l'image d'intérêt \(X\) est un vecteur aléatoire de
\(\mathbb{R}^d\) de densité \(f_X\) (quitte à réarranger les pixels et
les canaux RGB) et que le bruit \(Z\) est de densité \(f_Z\),
gaussienne, centrée, de covariance \(\Sigma\). L'image bruitée peut donc
être représentée par \(Y = X + Z\).

\runinhead{Question
1}\phantomsection\addcontentsline{toc}{subsubsection}{Question
1}\label{III-tweedie1}
Exprimer la densité \(f_Y\) de \(Y\) en fonction de \(f_X\) et
\(f_Z\).\solref{III-answer-tweedie1}

\runinhead{Question
2}\phantomsection\addcontentsline{toc}{subsubsection}{Question
2}\label{III-tweedie2}
Calculer le gradient de \(\log f_Y\), que l'on appelle aussi le score de
la loi de \(Y\).\solref{III-answer-tweedie2}

\runinhead{Question
3}\phantomsection\addcontentsline{toc}{subsubsection}{Question
3}\label{III-tweedie3}
En déduire une expression de \(\mathbb{E}(X|Y)\) ne faisant pas
intervenir la densité \(f_X\).

Cette propriété est particulièrement intéressante en pratique puisque
pour ce type de problèmes, on ne dispose généralement que d'obervations
bruitées qui vont nous permettre d'approximer les quantités nécessaires
pour effectuer ce calcul selon une approche statistique. Elle est par
ailleurs au coeur des algorithmes modernes de génération
d'image.\solref{III-answer-tweedie3}

\section{Non-réponse}\label{III-non-ruxe9ponse}

\emph{Inspiré du cours de probabilité de M. Christine (ENSAE
ParisTech).}

Un questionnaire est diffusé aux \(n\in\mathbb{N}^\ast\) étudiants de
l'école pour savoir combien de temps ils ont consacré à l'étude des
probabilités ce semestre. On note \(Y_i\) le temps de travail de
l'étudiant \(i \in \{1,\dots,n\}\) et \(X_i\) la variable valant \(1\)
s'il a répondu au questionnaire et \(0\) sinon. On suppose que les
\((X_1,Y_1),\dots,(X_n,Y_n)\) sont des vecteurs aléatoires indépendants
de même distribution qu'un vecteur générique \((X,Y)\) tel que

\begin{itemize}
\tightlist
\item
  \(X\) est une variable de Bernoulli de paramètre \(p\in\,]0,1[\)
  indiquant la probabilité de réponse,
\item
  \(Y\) est positive, de carré intégrable, d'espérance
  \(m\in\mathbb{R}_+\) et de variance \(\sigma^2 \in\mathbb{R}_+^\ast\).
  Le coefficient de corrélation entre \(X\) et \(Y\) est enfin noté
  \(\rho \in [-1,1]\).
\end{itemize}

\runinhead{Question
1}\phantomsection\addcontentsline{toc}{subsubsection}{Question
1}\label{III-nonrep1}
En reprenant la définition de l'espérance conditionnelle
\(\mathbb{E}(Y\mid X)\) comme meilleure approximation au sens des
moindres carrés de \(Y\) par une fonction de \(X\), montrer qu'elle
coïncide ici avec l'approximation affine de \(Y\) par \(X\) puis
l'écrire en fonction de \(m\), \(\rho\), \(\sigma\) et
\(p\).\solref{III-answer-nonrep1}

\runinhead{Question
2}\phantomsection\addcontentsline{toc}{subsubsection}{Question
2}\label{III-nonrep2}
On pose \(m_0 := \mathbb{E}(Y\mid X = 0)\) et
\(m_1 = \mathbb{E}(Y\mid X = 1)\). Calculer \(m_0\) et \(m_1\) en
fonction de \(m\), \(\rho\), \(\sigma\) et
\(p\).\solref{III-answer-nonrep2}

\runinhead{Question
3}\phantomsection\addcontentsline{toc}{subsubsection}{Question
3}\label{III-nonrep3}
On pose \(\sigma^2_0 := \mathbb{V}\left(Y\mid X = 0\right)\) et
\(\sigma^2_1 := \mathbb{V}\left(Y\mid X = 1\right)\). Vérifier l'égalité
\[\sigma^2 = \dfrac{(1-p)\,\sigma^2_0 + p\,\sigma^2_1}{1-\rho^2}.\]

\solref{III-answer-nonrep3}

\runinhead{Question
4}\phantomsection\addcontentsline{toc}{subsubsection}{Question
4}\label{III-nonrep4}
Que dire des résultats obtenus aux questions 2 et 3 lorsque~:

\begin{itemize}
\tightlist
\item
  \(X\) et \(Y\) sont non corrélées,
\item
  \(X\) et \(Y\) sont indépendantes\,?
\end{itemize}
\solref{III-answer-nonrep4}

\begin{center}\rule{0.5\linewidth}{0.5pt}\end{center}

